\problem{CX01}{随机快速排序的比较次数}{}
对 $n$ 个互异键使用随机快速排序：每个非空子问题独立均匀选择一个主元，与该子问题的其余每个键比较一次，再递归处理两侧。设总比较次数为 $C_n$。求 $\E C_n$ 的精确公式，并解释为什么平均为 $O(n\log n)$ 不代表每次运行都是这个量级。
\solution 按键值大小将键编号为 $1,\ldots,n$。令 $I_{ij}$ 表示键 $i<j$ 在算法中是否比较过。每次比较包含一个主元，该主元随后不再进入子问题，所以一对键至多比较一次，并且
\[
C_n=\sum_{1\le i<j\le n}I_{ij}.
\]
只观察秩区间 $\{i,i+1,\ldots,j\}$ 中最先被选作主元的键。区间外的主元不会将 $i,j$ 分到两侧；若第一个区间内主元是内部键，它立即分离 $i,j$，之后二者不再比较；若是端点，则两者当时比较。因此
\[
\Prob(I_{ij}=1)=\frac2{j-i+1}.
\]
“第一个均匀”可由等价实现解释：先给所有键赋独立连续随机优先级，每个子问题选择优先级最低者，所得主元选择规律相同；区间内最低优先级在各键间等可能。

线性性不需要各 $I_{ij}$ 独立。按间隔 $d=j-i$ 求和，令 $H_n=\sum_{k=1}^n1/k$，得到
\[
\begin{aligned}
\E C_n
&=2\sum_{d=1}^{n-1}\frac{n-d}{d+1}\\
&=2(n+1)(H_n-1)-2(n-1)\\
&=\boxed{2(n+1)H_n-4n}.
\end{aligned}
\]
公式对 $n=1$ 给出零；$n=0$ 单独定义 $C_0=0$。例如 $n=4$ 时均值为 $29/6$。

若每次主元恰为当前最小值，就会发生 $(n-1)+\cdots+1=n(n-1)/2$ 次比较；该事件具有正概率 $1/n!$。所以期望阶是对随机运行的平均陈述，不是逐次运行的最坏情形保证。
\insight 分析某一对键何时会相遇，比追踪整棵随机递归树更直接。两两指示量可以相关，期望仍可逐项相加。

\problem{CX02}{缺货与积压的最优平衡}{}
非负需求 $D$ 满足 $\E D<\infty$，其分布函数为 $F$。订购量 $q\ge0$ 可为实数，每单位缺货成本为 $c_u>0$，每单位积压成本为 $c_o>0$。最小化
\[
L(q)=c_u\E(D-q)_++c_o\E(q-D)_+.
\]
不假设需求分布连续，给出全部最优订购量。对于 $\Prob(D=0,1,2)=(1/4,1/2,1/4)$，分别求 $(c_u,c_o)=(3,1)$ 与 $(2,1)$ 时的解集。
\solution 每个固定需求下的成本都是 $q$ 的凸分段线性函数，因此 $L$ 也凸。左右差商均被常数 $\max(c_u,c_o)$ 控制，可以在期望内取单侧极限，得到
\[
\begin{aligned}
L'_-(q)&=(c_u+c_o)F(q-)-c_u,\\
L'_+(q)&=(c_u+c_o)F(q)-c_u.
\end{aligned}
\]
对 $q>0$，凸函数最优性的充要条件为左导数不正、右导数不负。因此令 $\tau=c_u/(c_u+c_o)$，全部最优量为
\[
\boxed{\{q\ge0:F(q-)\le\tau\le F(q)\}}.
\]
在边界 $q=0$，只需右导数不负；由于非负需求使 $F(0-)=0\le\tau$，同一公式依然适用。若 $q<0$，形式上延拓的成本随 $q$ 增加严格下降，所以不会遗漏受约束的最优点。

最优量存在：$\tau\in(0,1)$，广义分位数 $\inf\{x:F(x)\ge\tau\}$ 有限并满足上述条件。若分布函数恰在某段保持水平值 $\tau$，整个区间都可以最优；不能未经检查就宣布分位数唯一。

离散例中，$(c_u,c_o)=(3,1)$ 给出 $\tau=3/4$。在 $q=1$ 处 $F$ 跳到 $3/4$，到 $2$ 前维持此值，故全部最优量为 $[1,2]$；两端及其间的期望成本均为 $1$。当 $(c_u,c_o)=(2,1)$ 时 $\tau=2/3$ 严格落在 $q=1$ 的跳跃内，故唯一解为 $q=1$，最小成本为 $3/4$。
\insight 非对称绝对损失选出成本加权分位数。需求均值一般不是这个决策问题的最优答案，原子和分布空档还会影响唯一性。

\problem{CX03}{给定均值的最大熵分布}{}
随机变量 $N$ 取非负整数，且 $\E N=\mu>0$。用自然对数定义熵
\[
H(N)=-\sum_{k=0}^\infty p_k\log p_k,\qquad 0\log0=0.
\]
证明熵有有限上界，求达到最大熵的唯一分布及最大值。允许使用有限分布的对数和不等式，但须处理无限支撑，不得预先假设熵有限。
\solution 令 $r=\mu/(1+\mu)\in(0,1)$，考虑几何概率
\[
q_k=(1-r)r^k,\qquad k\ge0.
\]
它的均值为 $r/(1-r)=\mu$。若 $N$ 服从该分布，则
\[
H(q)=-\log(1-r)-\mu\log r
=(1+\mu)\log(1+\mu)-\mu\log\mu.
\]

为在一般分布下严格给出上界，将状态 $0,\ldots,M-1$ 分开，把 $N\ge M$ 合并成一个尾部状态，记尾概率为 $t_M$。参考分布 $q$ 对该尾部赋概率 $r^M$。有限分布的对数和不等式给出分组熵
\[
H_M\le-\sum_{k<M}p_k\log q_k-t_M\log(r^M).
\]
右端等于
\[
-(1-t_M)\log(1-r)
-\log r\left(\sum_{k<M}kp_k+Mt_M\right).
\]
当 $M\to\infty$，有 $t_M\to0$，而括号内是 $\E\min(N,M)\uparrow\mu$。同时
$H_M=-\sum_{k<M}p_k\log p_k-t_M\log t_M$ 趋于 $H(N)$，即使起初允许后者为无穷也成立，因为 $-t\log t\to0$。取极限得
\[
H(N)\le(1+\mu)\log(1+\mu)-\mu\log\mu<\infty.
\]

还要证明唯一性。以上右端与 $H_M$ 的差是分组相对熵 $D_M\ge0$，其极限等于 $H(q)-H(N)$。若 $p\ne q$，选一个 $j$ 使 $p_j\ne q_j$。对每个 $M>j$，把除 $j$ 外的所有分组再合并，对数和不等式给出
\[
D_M\ge
p_j\log\frac{p_j}{q_j}
+(1-p_j)\log\frac{1-p_j}{1-q_j}>0.
\]
该严格正下界不依赖 $M$，故不可能取等。于是几何分布是唯一最大熵分布。
\insight 在无限状态上，先将尾部合并成一个状态，再取极限，可避免把待证的有限熵假设偷偷用于相对熵运算。
